\chapter{Poisson Processes}

So far almost every random object in this book has lived at a single moment.
A coin flip, a normal measurement, a Poisson count of typos on a page --- each
is one snapshot. But the world does not hand us snapshots. It hands us
\emph{streams}: customers walking through a door, photons striking a detector,
earthquakes rattling a fault, requests hitting a server, mutations accumulating
along a strand of DNA. Events arrive, one after another, scattered along the
axis of time, and we want a probability model for the whole pattern at once.

The \emph{Poisson process} is that model. It is the most important random process
in applied probability --- the ``hydrogen atom'' of stochastic processes. It is
what you get when events happen completely at random in time, with no memory and
no coordination, at some steady average rate.
It is simple enough to compute almost everything by hand, yet rich enough to be
the foundation of queueing theory, reliability engineering, insurance
mathematics, particle physics, and the analysis of point patterns in space.

This chapter builds the Poisson process three times over, from three different
starting points, and then proves that all three describe the \emph{same} object.
That triple viewpoint is not pedantic repetition --- each lens makes a different
property obvious, and the deepest understanding comes from being able to slide
between them at will.

\section{Three ways to see the same thing}
\label{ch25:three-views}

Picture a long horizontal time axis and a scattering of dots on it --- the
instants at which events occur. There are three natural questions you can ask
about that picture, and each one is a complete description of it.

\begin{enumerate}
\item \textbf{The counting view.} Fix a window of time, say the interval
$[0,t]$, and ask: \emph{how many} dots fall inside it? Call that number $N(t)$.
As $t$ grows, $N(t)$ ticks upward in unit jumps. The whole process is the random
staircase $t \mapsto N(t)$.
\item \textbf{The interarrival view.} Forget windows. Just read off the
\emph{gaps}: how long until the first dot, then how long from the first to the
second, and so on. Call these waiting times $T_1, T_2, T_3, \dots$. Knowing all
the gaps is the same as knowing where every dot is.
\item \textbf{The infinitesimal view.} Zoom in on a tiny sliver of time of width
$\dd t$. Ask only one local question: \emph{what is the chance a dot lands in
this sliver}, and does what happens here depend on what happened elsewhere? This
view describes the process by its instantaneous behavior, like a differential
equation.
\end{enumerate}

\begin{intuitionbox}
These are the same information packaged three ways, exactly as a function can be
described by its values, its inverse, or its derivative. The counting view is
global and discrete (a staircase). The interarrival view is a list of gaps. The
infinitesimal view is local and ``differential.'' The Poisson process is the one
process on which all three descriptions are as clean and as friendly as
possible, and the central point of the chapter is that cleanness in any one view
\emph{forces} cleanness in the other two.
\end{intuitionbox}

\begin{figure}[H]
\centering
\begin{tikzpicture}[scale=1]
% time axis
\draw[->,thick] (0,0) -- (11,0) node[right]{time $t$};
% event dots (arrival instants)
\foreach \x in {1.1, 1.7, 3.0, 3.4, 3.7, 5.5, 7.1, 8.0, 9.4, 9.8} {
  \filldraw[densityred] (\x,0) circle (2.3pt);
  \draw[densityred, thin] (\x,0) -- (\x,-0.18);
}
% interarrival braces for first few
\draw[decorate,decoration={brace,amplitude=5pt},sourceblue]
  (0,0.35) -- (1.1,0.35) node[midway,above=4pt,sourceblue,font=\small]{$T_1$};
\draw[decorate,decoration={brace,amplitude=5pt},sourceblue]
  (1.1,0.35) -- (1.7,0.35) node[midway,above=4pt,sourceblue,font=\small]{$T_2$};
\draw[decorate,decoration={brace,amplitude=5pt},sourceblue]
  (1.7,0.35) -- (3.0,0.35) node[midway,above=4pt,sourceblue,font=\small]{$T_3$};
\node[densityred,font=\small] at (5.5,-0.45) {arrival instants};
\end{tikzpicture}
\caption{A single realization of a point process on the time line. The counting
view asks how many red dots lie left of $t$; the interarrival view records the
blue gaps $T_1, T_2, \dots$; the infinitesimal view asks about a tiny window
around one point.}
\label{ch25:fig-pointprocess}
\end{figure}

\section{The infinitesimal definition: events at random, with no memory}
\label{ch25:infinitesimal}

We start with the infinitesimal view because it captures most directly what we
\emph{mean} by ``events happening at random at a steady rate.'' Three
assumptions, each a plain-language statement about local behavior, pin the whole
process down.

\begin{definition}[Poisson process, infinitesimal form]
\label{ch25:def-infinitesimal}
A counting process $\{N(t): t \ge 0\}$ with $N(0)=0$ is a \emph{(homogeneous)
Poisson process with rate} $\lambda > 0$ if, for a small time increment $h$,
\begin{enumerate}
\item \textbf{(one at a time)} the chance of exactly one event in a window of
width $h$ is $\lambda h + o(h)$;
\item \textbf{(no doubles)} the chance of two or more events in a window of width
$h$ is $o(h)$ --- vanishingly small compared to $h$;
\item \textbf{(independence)} what happens in disjoint time windows is
independent, and the rule above does not depend on \emph{where} the window sits
(stationarity).
\end{enumerate}
Here $o(h)$ denotes a quantity that goes to zero faster than $h$, i.e.
$o(h)/h \to 0$ as $h \to 0$.
\end{definition}

Read those three lines as a story. Events come \emph{singly} --- two of them
never share an instant. They come at a \emph{constant intensity} $\lambda$: in a
window of width $h$ you expect about $\lambda h$ of them, so $\lambda$ is events
per unit time. And the future is \emph{indifferent to the past}: knowing the
process has been quiet for an hour tells you nothing about the next second.

\begin{happeningbox}
The condition $\P(\text{one event in }h) = \lambda h + o(h)$ is the engine of the
entire chapter. It says the probability of an arrival is, to first order,
\emph{proportional to the length of time you wait} --- and proportional to
nothing else. Not to how long since the last arrival, not to the clock reading,
not to how many have already occurred. ``Random'' here means precisely this:
each instant gets the same, independent, tiny shot at producing an event. Every
clean formula below is just this local rule integrated up.
\end{happeningbox}

\subsection{From the local rule to the count}

We can extract the distribution of $N(t)$ directly from
Definition~\ref{ch25:def-infinitesimal}. Chop $[0,t]$ into $n$ equal slices of
width $h = t/n$. By the ``no doubles'' rule, when $n$ is enormous each slice
holds either zero events or one, and the chance of one is $\lambda h + o(h)
\approx \lambda t / n$. By independence, the slices behave like $n$ independent
coin flips, each landing ``event'' with probability $p = \lambda t/n$. So
\[
N(t) \;\approx\; \Bin\!\left(n,\ \tfrac{\lambda t}{n}\right).
\]
Now let $n \to \infty$. This is exactly the limit that \emph{defines} the Poisson
distribution --- a binomial with many trials and a tiny per-trial probability,
holding the mean $np = \lambda t$ fixed. The law of small numbers gives
\[
\P\big(N(t) = k\big) \;=\; \lim_{n\to\infty}\binom{n}{k}
\left(\tfrac{\lambda t}{n}\right)^{k}\!\left(1-\tfrac{\lambda t}{n}\right)^{n-k}
\;=\; e^{-\lambda t}\,\frac{(\lambda t)^{k}}{k!}.
\]

\begin{theorem}[Counts are Poisson]
\label{ch25:thm-counts-poisson}
In a rate-$\lambda$ Poisson process, the number of events in any window of length
$t$ is
\[
N(t) \sim \Pois(\lambda t), \qquad
\P\big(N(t)=k\big) = e^{-\lambda t}\frac{(\lambda t)^k}{k!}, \quad k=0,1,2,\dots
\]
More generally, the number of events in any interval $(a,b]$ is
$\Pois\big(\lambda(b-a)\big)$, and counts in disjoint intervals are independent.
\end{theorem}

\begin{whybox}
Why does the slicing argument become exact in the limit rather than merely
approximate? Two reasons, and both are the $o(h)$ terms doing their job. First,
the probability that any slice contains \emph{two or more} events is $o(h)$ per
slice, so across all $n$ slices the total chance of a ``double'' is at most
$n\cdot o(t/n) = t\cdot \frac{o(t/n)}{t/n} \to 0$. The single-event approximation
is therefore not lossy in the limit. Second, the $o(h)$ correction to the
single-event probability likewise washes out. What survives is precisely a
$\Bin(n, \lambda t/n)$, and its limit is Poisson. The Poisson distribution is not
\emph{assumed} here --- it is \emph{forced} by ``rare, independent, one-at-a-time,
steady-rate.''
\end{whybox}

This is why the distribution and the process share a name. The Poisson
\emph{distribution} of Chapter~7 counts rare independent events; the Poisson
\emph{process} is what you get when you let those events live on a time line and
ask about every window at once. The number $\lambda t$ is the expected count in a
window of length $t$, and $\E N(t) = \Var N(t) = \lambda t$ --- the count's mean
and variance both equal $\lambda t$, the signature of a Poisson variable
(Figure~\ref{ch25:fig-poisson-pmf}).

\begin{formalbox}
\textbf{The differential equation derivation (Kolmogorov forward equation).}
We can also derive the Poisson pmf directly from calculus without taking binomial limits. Let $P_n(t) = \P(N(t) = n)$ be the probability of observing exactly $n$ events by time $t$. Over a tiny interval $[t, t+h]$, independent increments and the infinitesimal definition give:
\[
  P_0(t+h) = P_0(t)\,\P(0 \text{ in } h) = P_0(t)(1 - \lambda h + o(h)).
\]
Rearranging and taking the limit $h \to 0$ yields the differential equation:
\[
  P_0'(t) = -\lambda P_0(t) \quad \Longrightarrow \quad P_0(t) = e^{-\lambda t} \quad (\text{using } P_0(0) = 1).
\]
For $n \ge 1$, reaching state $n$ at time $t+h$ occurs either by being at $n$ and seeing $0$ arrivals, or being at $n-1$ and seeing $1$ arrival (all other transition paths have probability $o(h)$):
\[
  P_n(t+h) = P_n(t)(1 - \lambda h) + P_{n-1}(t)(\lambda h) + o(h).
\]
Subtracting $P_n(t)$, dividing by $h$, and taking $h \to 0$ gives:
\[
  P_n'(t) + \lambda P_n(t) = \lambda P_{n-1}(t).
\]
Multiplying through by the integrating factor $e^{\lambda t}$:
\[
  \frac{\dd}{\dd t}\big[ e^{\lambda t} P_n(t) \big] = \lambda e^{\lambda t} P_{n-1}(t).
\]
Since $P_0(t) = e^{-\lambda t}$, for $n=1$ we have $\frac{\dd}{\dd t}[e^{\lambda t} P_1(t)] = \lambda$, which integrates with initial condition $P_1(0) = 0$ to $e^{\lambda t} P_1(t) = \lambda t$, so $P_1(t) = (\lambda t) e^{-\lambda t}$. By induction:
\[
  P_n(t) = \frac{(\lambda t)^n}{n!} e^{-\lambda t}.
\]
This system of differential equations is the prototype for the generator equations of continuous-time Markov chains in Chapter~27.
\end{formalbox}

\begin{figure}[H]
\centering
\begin{tikzpicture}
\begin{axis}[
  width=12cm, height=5.2cm,
  ybar, bar width=10pt,
  xmin=-0.6, xmax=12.6, ymin=0, ymax=0.24,
  xtick={0,2,4,6,8,10,12},
  xlabel={number of events $k$ in the window},
  ylabel={$\P(N(t)=k)$},
  title={The count is Poisson: pmf of $N(t)\sim\Pois(\lambda t)$ at mean $\lambda t=4$},
  axis lines=left, title style={font=\small},
]
% Poisson(4) pmf: e^-4 4^k / k!
\addplot[fill=histfill, draw=distblue] coordinates {
 (0,0.018316) (1,0.073263) (2,0.146525) (3,0.195367) (4,0.195367)
 (5,0.156293) (6,0.104196) (7,0.059540) (8,0.029770) (9,0.013231)
 (10,0.005292) (11,0.001925) (12,0.000642)};
\end{axis}
\end{tikzpicture}
\caption{The number of events in a window of length $t$ follows the Poisson pmf
with mean $\lambda t$ (here $\lambda t = 4$). The distribution peaks near its
mean and has the equal mean-equals-variance signature: both equal $\lambda t$.}
\label{ch25:fig-poisson-pmf}
\end{figure}

\begin{figure}[H]
\centering
\begin{tikzpicture}
\begin{axis}[
  width=12cm, height=5.6cm,
  ymin=0, ymax=10.5, xmin=0, xmax=10,
  xlabel={time $t$}, ylabel={$N(t)$},
  title={A sample path of the counting process (a random staircase)},
  axis lines=left, ytick={0,2,4,6,8,10},
  title style={font=\small},
]
% staircase: jumps at the arrival instants from fig above
\addplot[const plot, no marks, thick, sourceblue] coordinates {
 (0,0) (1.1,1) (1.7,2) (3.0,3) (3.4,4) (3.7,5) (5.5,6) (7.1,7) (8.0,8)
 (9.4,9) (9.8,10) (10,10)};
% mark arrival instants on the x-axis
\addplot[only marks, mark=*, mark size=1.6pt, densityred] coordinates {
 (1.1,0)(1.7,0)(3.0,0)(3.4,0)(3.7,0)(5.5,0)(7.1,0)(8.0,0)(9.4,0)(9.8,0)};
\end{axis}
\end{tikzpicture}
\caption{The counting view of the same realization as
Figure~\ref{ch25:fig-pointprocess}: $N(t)$ jumps up by one at each red arrival
and is flat in between. The Poisson process \emph{is} this staircase.}
\label{ch25:fig-staircase}
\end{figure}

\section{The interarrival view: why the gaps are exponential}
\label{ch25:interarrival}

Now switch lenses. Instead of counting, watch the clock between events. Let $T_1$
be the time of the first event, $T_2$ the gap from the first to the second, and
so on. What these gaps turn out to be is the most striking fact about the Poisson
process.

\begin{theorem}[Interarrival times are i.i.d.\ exponential]
\label{ch25:thm-exponential}
In a rate-$\lambda$ Poisson process the interarrival times
$T_1, T_2, T_3, \dots$ are independent and identically distributed, each
$\Expo(\lambda)$:
\[
\P(T_i > s) = e^{-\lambda s}, \qquad f_{T_i}(s) = \lambda e^{-\lambda s},
\quad s \ge 0,
\]
with mean gap $\E T_i = 1/\lambda$.
\end{theorem}

\begin{whybox}
The proof of the first gap is a one-liner once you have the counting view. The
event ``the first arrival happens after time $s$'' is \emph{exactly} the event
``no arrivals in $[0,s]$.'' But the count $N(s)$ is $\Pois(\lambda s)$, so
\[
\P(T_1 > s) = \P\big(N(s) = 0\big) = e^{-\lambda s}\frac{(\lambda s)^0}{0!}
= e^{-\lambda s}.
\]
That is the exponential survival function on the nose. So $T_1 \sim
\Expo(\lambda)$, and differentiating gives the density $\lambda e^{-\lambda s}$.
For the later gaps, independence of disjoint intervals means that once the first
event lands, the process from that instant onward is a fresh Poisson process that
has forgotten its past entirely --- so $T_2$ is, by the identical argument,
$\Expo(\lambda)$ and independent of $T_1$, and so on for all the gaps.
\end{whybox}

Notice that the exponential distribution did not get assumed; it got
\emph{derived} from ``no arrivals yet'' plus ``counts are Poisson.'' And the
reverse holds too: if you simply \emph{define} a process by dropping i.i.d.\
$\Expo(\lambda)$ gaps end to end, you recover a process whose counts are
$\Pois(\lambda t)$. This is the equivalence promised in the introduction.

\begin{intuitionbox}
The deep reason the gaps are exponential is \emph{memorylessness}. The
infinitesimal rule says every instant is an independent fresh shot at an event,
indifferent to history. A waiting time built from a memoryless local rule must
itself be memoryless: $\P(T > s+u \given T > s) = \P(T > u)$. And the exponential
is the \emph{only} continuous distribution with that property (Chapter~11). So
``random events with no memory'' and ``exponential gaps'' are two names for the
same thing.
\end{intuitionbox}

\subsection{Where the $n$-th event lands: the Erlang/Gamma}

The time of the $n$-th arrival is the sum of the first $n$ gaps,
\[
S_n = T_1 + T_2 + \cdots + T_n.
\]
A sum of $n$ independent $\Expo(\lambda)$ variables is, by the convolution
results of Chapter~17, a Gamma with shape $n$ --- often called the \emph{Erlang}
distribution:
\[
S_n \sim \Gam(n, \lambda), \qquad
f_{S_n}(t) = \frac{\lambda^n t^{\,n-1} e^{-\lambda t}}{(n-1)!}, \quad t \ge 0.
\]
There is a consistency check hiding here, the bridge between the interarrival and
counting views: the $n$-th event lands after time $t$ exactly when fewer than $n$
events have occurred by time $t$. In symbols,
$\{S_n > t\} = \{N(t) < n\}$, so
\[
\P(S_n > t) = \P\big(N(t) \le n-1\big)
= \sum_{k=0}^{n-1} e^{-\lambda t}\frac{(\lambda t)^k}{k!}.
\]
Differentiating that finite sum collapses (the series telescopes) right back to
the Erlang density above --- the two viewpoints agreeing to the last term.

\begin{figure}[H]
\centering
\begin{tikzpicture}
\begin{axis}[
  width=12cm, height=5.4cm,
  domain=0:12, samples=160,
  xmin=0, xmax=12, ymin=0, ymax=0.42,
  xlabel={$t$}, ylabel={density of $S_n$},
  title={Arrival-time densities $S_n\sim\Gam(n,1)$ for $n=1,2,3,5$},
  axis lines=left, legend pos=north east,
  title style={font=\small}, legend style={font=\small},
]
\addplot[densityred, very thick] {exp(-x)};
\addlegendentry{$n=1$ (exponential)}
\addplot[distblue, very thick] {x*exp(-x)};
\addlegendentry{$n=2$}
\addplot[targetgreen, very thick] {x^2*exp(-x)/2};
\addlegendentry{$n=3$}
\addplot[quotientorange, very thick] {x^4*exp(-x)/24};
\addlegendentry{$n=5$}
\end{axis}
\end{tikzpicture}
\caption{The time of the $n$-th arrival is a sum of $n$ exponential gaps, so its
density is Erlang/Gamma. As $n$ grows the hump marches right (the $n$-th event
takes longer) and grows more symmetric --- a preview of the central limit theorem
acting on the gaps.}
\label{ch25:fig-erlang}
\end{figure}

\begin{example}[Phone calls at a help desk]
\label{ch25:ex-helpdesk}
Calls arrive as a Poisson process at $\lambda = 6$ per hour. (a) What is the
chance of no calls in a given 10-minute stretch? (b) What is the expected wait
until the next call? (c) What is the chance the third call of the day arrives
within the first half hour?

\emph{Solution.} Work in hours. (a) A 10-minute window is $t = 1/6$ hour, so the
count is $\Pois(6 \cdot \tfrac16) = \Pois(1)$ and $\P(N=0) = e^{-1} \approx
0.368$. (b) The gap is $\Expo(6)$ with mean $1/6$ hour $= 10$ minutes. (c) We
want $\P(S_3 \le 0.5)$ where $S_3 \sim \Gam(3,6)$. Equivalently, by the bridge
identity, we want $\P\big(N(0.5) \ge 3\big)$ with $N(0.5) \sim \Pois(3)$:
\[
\P\big(N(0.5)\ge 3\big) = 1 - e^{-3}\Big(1 + 3 + \tfrac{9}{2}\Big)
= 1 - 8.5\,e^{-3} \approx 0.577.
\]
Notice how the counting view turned a Gamma integral into a three-term sum ---
that switch between lenses is the practical payoff of the equivalence.
\end{example}

\begin{historybox}
The exponential-gap picture has deep engineering roots. Around 1909 the Danish
mathematician and telephone engineer \emph{Agner Krarup Erlang}, working for the
Copenhagen Telephone Company, modeled incoming calls as random arrivals to size
the number of circuits a switchboard needed. His insight --- that call arrivals
have no memory and that waiting times are exponential --- launched the entire
field of \emph{queueing theory}. The Erlang distribution and the unit of
telephone traffic (the \emph{erlang}) both carry his name. The same mathematics
that once dimensioned telephone exchanges now sizes data-center capacity, call
centers, and hospital emergency rooms.
\end{historybox}

\section{The two structural superpowers: superposition and thinning}
\label{ch25:super-thin}

Two operations make the Poisson process behave like a perfect Lego brick: you can
\emph{merge} processes and \emph{split} a process, and in both cases you land back
on Poisson processes. These are not curiosities --- they are the reason the model
appears everywhere, because real streams are constantly being pooled and
filtered.

\begin{theorem}[Superposition]
\label{ch25:thm-superposition}
If $N_1, N_2, \dots, N_m$ are independent Poisson processes with rates
$\lambda_1, \dots, \lambda_m$, then their pooled process $N = N_1 + \cdots + N_m$
--- the single stream you get by ignoring labels and watching all events together
--- is a Poisson process with rate $\lambda_1 + \cdots + \lambda_m$.
\end{theorem}

The infinitesimal view makes this almost obvious. In a tiny window, the chance of
an event from source $i$ is $\lambda_i h + o(h)$; the chance of an event from
\emph{some} source is the sum $(\sum_i \lambda_i)h + o(h)$ (the chance two
different sources fire in the same sliver is $o(h)$, negligible). Independence of
disjoint windows is inherited from the parts. All three axioms of
Definition~\ref{ch25:def-infinitesimal} hold for the pool with the summed rate.
Rates simply add.

\begin{figure}[H]
\centering
\begin{tikzpicture}[
  scale=1,
  lane/.style={draw=gray!30, fill=gray!3, rounded corners=3pt},
  ruler/.style={font=\scriptsize, gray}
]
  % Ruler on top
  \draw[->, thick, gray] (2.2, 3.4) -- (11.8, 3.4) node[right, font=\scriptsize] {time $t$};
  \foreach \t in {0, 2, 4, 6, 8, 10} {
    \draw[gray] (2.2 + \t*0.9, 3.35) -- (2.2 + \t*0.9, 3.45);
    \node[ruler, above] at (2.2 + \t*0.9, 3.45) {\t};
  }

  % Stream 1 (top lane)
  \node[font=\footnotesize\bfseries, sourceblue, anchor=east] at (2.0, 2.6) {Stream 1 ($N_1$)};
  \node[font=\scriptsize, sourceblue, anchor=east] at (2.0, 2.25) {rate $\lambda_1 = 0.8$};
  \draw[thick, sourceblue] (2.2, 2.4) -- (11.6, 2.4);
  \foreach \x in {0.8, 2.1, 4.3, 6.2, 7.9, 9.8} {
    \filldraw[sourceblue] (2.2 + \x*0.9, 2.4) circle (2.2pt);
    \draw[sourceblue, dashed, opacity=0.4] (2.2 + \x*0.9, 2.4) -- (2.2 + \x*0.9, 0.4);
  }

  % Stream 2 (middle lane)
  \node[font=\footnotesize\bfseries, quotientorange, anchor=east] at (2.0, 1.5) {Stream 2 ($N_2$)};
  \node[font=\scriptsize, quotientorange, anchor=east] at (2.0, 1.15) {rate $\lambda_2 = 0.5$};
  \draw[thick, quotientorange] (2.2, 1.3) -- (11.6, 1.3);
  \foreach \x in {1.4, 3.5, 5.4, 8.7} {
    \filldraw[quotientorange] (2.2 + \x*0.9, 1.3) circle (2.2pt);
    \draw[quotientorange, dashed, opacity=0.4] (2.2 + \x*0.9, 1.3) -- (2.2 + \x*0.9, 0.4);
  }

  % Merged Stream (bottom lane)
  \node[font=\footnotesize\bfseries, densityred, anchor=east] at (2.0, 0.5) {Superposition ($N$)};
  \node[font=\scriptsize, densityred, anchor=east] at (2.0, 0.15) {rate $\lambda = \lambda_1 + \lambda_2 = 1.3$};
  \draw[thick, densityred] (2.2, 0.4) -- (11.6, 0.4);
  \foreach \x in {0.8, 2.1, 4.3, 6.2, 7.9, 9.8} {
    \filldraw[sourceblue] (2.2 + \x*0.9, 0.4) circle (2.2pt);
  }
  \foreach \x in {1.4, 3.5, 5.4, 8.7} {
    \filldraw[quotientorange] (2.2 + \x*0.9, 0.4) circle (2.2pt);
  }

\end{tikzpicture}
\caption{Superposition of two independent Poisson processes. Stream 1 (blue dots, rate $\lambda_1$) and Stream 2 (orange dots, rate $\lambda_2$) merge into a single pooled stream (bottom). The combined arrivals form a Poisson process with rate $\lambda_1 + \lambda_2$. The probability that two events collide at the exact same instant is zero ($o(h)$ infinitesimal double-arrival exclusion).}
\label{ch25:fig-superposition}
\end{figure}

\begin{theorem}[Thinning]
\label{ch25:thm-thinning}
Take a rate-$\lambda$ Poisson process and independently label each event ``type
A'' with probability $p$ and ``type B'' with probability $1-p$, the coin flips
independent of everything else. Then the type-A events form a Poisson process of
rate $\lambda p$, the type-B events form a Poisson process of rate
$\lambda(1-p)$, and --- the surprising part --- \emph{the two thinned processes
are independent of each other}.
\end{theorem}

\begin{whybox}
Why are the two thinned streams independent, when they came from a single shared
process? Condition on the total count $N(t)=n$. Given $n$ events, each is
independently type A with probability $p$, so the number of type-A events is
$\Bin(n,p)$. Summing over $n$ against the Poisson law of $N(t)$:
\[
\P(N_A = j,\ N_B = k)
= \underbrace{e^{-\lambda t}\frac{(\lambda t)^{j+k}}{(j+k)!}}_{\P(N=j+k)}\cdot
\underbrace{\binom{j+k}{j}p^j(1-p)^k}_{\text{labelling}}.
\]
Now do the algebra: $(j+k)!$ cancels, $(\lambda t)^{j+k}=(\lambda t)^j(\lambda
t)^k$ splits, and $e^{-\lambda t}=e^{-\lambda p t}e^{-\lambda(1-p)t}$ factors. The
joint probability separates into
\[
\Big[e^{-\lambda p t}\tfrac{(\lambda p t)^j}{j!}\Big]
\cdot
\Big[e^{-\lambda (1-p) t}\tfrac{(\lambda (1-p) t)^k}{k!}\Big],
\]
a product of two independent Poisson pmfs with means $\lambda p t$ and
$\lambda(1-p)t$. The factorization \emph{is} the independence. This ``Poissonizing
the trials'' trick --- a random Poisson number of binomial labels splits into
independent Poissons --- is one of the more surprising identities in probability.
\end{whybox}

\begin{figure}[H]
\centering
\begin{tikzpicture}[scale=1, every node/.style={font=\small}]
% pooled / superposition (top)
\draw[->,thick] (0,2.2) -- (7.2,2.2);
\foreach \x in {0.6,1.4,2.1,3.3,4.0,5.1,6.0,6.6} \filldraw[densityred] (\x,2.2) circle (2pt);
\node[densityred] at (8.4,2.2) {rate $\lambda$};
% thinning arrows
\draw[->,thick,gray] (3.6,1.95) -- (3.6,1.25);
\node[gray] at (5.7,1.55) {label A w.p. $p$, B w.p. $1-p$};
% type A
\draw[->,thick] (0,0.6) -- (7.2,0.6);
\foreach \x in {0.6,2.1,4.0,6.0} \filldraw[distblue] (\x,0.6) circle (2pt);
\node[distblue] at (8.6,0.6) {A: rate $\lambda p$};
% type B
\draw[->,thick] (0,-0.6) -- (7.2,-0.6);
\foreach \x in {1.4,3.3,5.1,6.6} \filldraw[targetgreen] (\x,-0.6) circle (2pt);
\node[targetgreen] at (8.8,-0.6) {B: rate $\lambda(1{-}p)$};
\end{tikzpicture}
\caption{Thinning: each event of a rate-$\lambda$ process is independently
sorted into stream A or stream B. The two sub-streams are themselves independent
Poisson processes whose rates split the original. Run the picture upward and it
reads as superposition: two independent Poisson streams merge into one whose rate
is the sum.}
\label{ch25:fig-thinning}
\end{figure}

\begin{example}[Splitting a traffic stream]
\label{ch25:ex-thinning}
Cars pass a sensor as a Poisson process at $\lambda = 120$ per hour; each
independently turns left with probability $0.25$. Then left-turners form a
Poisson process at $0.25 \cdot 120 = 30$ per hour and straight-throughs at $90$
per hour, and these two streams are independent. So the number of left-turners in
the next hour is $\Pois(30)$, completely uninformative about the number going
straight. The pooled-then-split structure means you can model each lane in
isolation.
\end{example}

\begin{applicationbox}[Modeling rare events, arrivals, and queues]
The Poisson process is the default first model whenever many independent sources
each rarely contribute to a common stream --- precisely the setting of the law of
small numbers.
\begin{itemize}
\item \textbf{Queues and servers.} An ``M/M/1'' queue (the workhorse of
performance modeling) has \emph{M}arkovian Poisson arrivals and \emph{M}arkovian
exponential service. Web-request load, database transactions, and packets on a
link are all routinely modeled this way; Chapter~27 develops the queueing theory
that sits on top of this chapter.
\item \textbf{Reliability and insurance.} Failures of a large fleet of
independent components, or claims arriving at an insurer, are classic thinned and
superposed Poisson streams; the total claim amount over a year is a
\emph{compound Poisson} sum.
\item \textbf{Physics and biology.} Radioactive decay clicks on a Geiger counter,
photons hitting a detector, and mutations along a genome are textbook Poisson
processes --- the independence of disjoint windows reflects genuinely
memoryless physics.
\item \textbf{Networks and ML.} Arrival models for recommendation impressions, ad
auctions, and event logs start from Poisson assumptions; Hawkes processes (where
events \emph{excite} more events) are the standard generalization when the
independence assumption breaks.
\end{itemize}
\end{applicationbox}

\section{The conditional-uniformity property}
\label{ch25:uniformity}

The next fact surprises almost everyone the first time, and once it sinks in it
demystifies the whole process.

\begin{theorem}[Order statistics / conditional uniformity]
\label{ch25:thm-uniformity}
Condition on $N(t)=n$: that is, suppose you are told exactly $n$ events occurred
in $[0,t]$, but not when. Then the $n$ arrival times, regarded as an unordered
set, are distributed exactly like $n$ independent $\Unif(0,t)$ points. Their
sorted values $S_1 < \cdots < S_n$ are the order statistics of $n$ uniform draws.
\end{theorem}

In words: \emph{given the count}, the Poisson process scatters its events
completely uniformly and with no preference for any sub-interval. The rate
$\lambda$ governs \emph{how many} events occur; once that number is fixed,
$\lambda$ drops out entirely and the events are just uniform confetti.

\begin{whybox}
A clean way to see it: chop $[0,t]$ into the slices of the infinitesimal view.
Each event is equally likely to fall in any slice (the rate is constant), and the
events are placed independently. ``Equally likely across the interval,
independent'' is the definition of i.i.d.\ uniform. More formally, write the joint
density of getting $n$ events near specified ordered times $s_1<\cdots<s_n$ and no
events elsewhere; the rate factors $\lambda^n$ and the empty-interval
probabilities multiply to $e^{-\lambda t}$, and after dividing by
$\P(N(t)=n)=e^{-\lambda t}(\lambda t)^n/n!$ everything $\lambda$-dependent
cancels, leaving the constant density $n!/t^n$ --- exactly the joint density of
$n$ ordered uniforms.
\end{whybox}

\begin{paradoxbox}[The inspection / waiting-time paradox]
You arrive at a bus stop at a random moment. Buses come as a Poisson process with
mean gap $1/\lambda$. How long is the gap of time that \emph{you} land inside?
Naively you might say $1/\lambda$, since that is the average gap. But the answer
is $2/\lambda$ --- on average the bus you wait for took \emph{twice} the mean gap
to arrive.

How can the gap you land in be longer than the average gap? Because your random
arrival is more likely to fall inside a \emph{long} gap than a short one --- a
long gap is, by definition, a bigger target. This \emph{length-biased sampling}
inflates the gap you observe. Concretely, by memorylessness the time you still
wait for the next bus is $\Expo(\lambda)$ with mean $1/\lambda$; but the time
since the \emph{previous} bus is \emph{also} $\Expo(\lambda)$ with mean
$1/\lambda$ (the process is memoryless looking backward too). The gap containing
you is the sum, mean $2/\lambda$. The lesson generalizes far beyond buses: when
you sample by ``whatever interval I happen to be in,'' you oversample the long
ones. It is the same bias that makes your friends seem to have more friends than
you do, and that makes class sizes feel larger than the average roster
(Figure~\ref{ch25:fig-inspection}).
\end{paradoxbox}

\begin{figure}[H]
\centering
\begin{tikzpicture}[scale=1, every node/.style={font=\small}]
% time axis with bus arrivals making gaps of varying length
\draw[->,thick] (0,0) -- (11,0) node[right]{time};
\foreach \x in {0.5, 1.3, 2.0, 5.5, 6.2, 7.0, 7.6, 10.2} {
  \filldraw[densityred] (\x,0) circle (2.2pt);
  \draw[densityred] (\x,0) -- (\x,-0.15);
}
% highlight the long gap 2.0 -- 5.5
\draw[decorate,decoration={brace,amplitude=5pt,mirror},structurepurple,thick]
  (2.0,-0.35) -- (5.5,-0.35) node[midway,below=4pt,structurepurple]{long gap};
% the random observer arrow lands in the long gap
\draw[->,thick,targetgreen] (3.7,1.0) -- (3.7,0.12);
\node[targetgreen] at (3.7,1.25) {you arrive here};
\node[densityred] at (6.5,0.55) {bus arrivals};
\end{tikzpicture}
\caption{Length-biased sampling behind the inspection paradox. A uniformly
random arrival time is more likely to land inside a \emph{wide} gap than a narrow
one, because wide gaps cover more of the axis. So the gap you observe is, on
average, longer than a gap picked at random from the list --- mean $2/\lambda$
rather than $1/\lambda$.}
\label{ch25:fig-inspection}
\end{figure}

\section{Nonhomogeneous Poisson processes: a rate that breathes}
\label{ch25:nonhomogeneous}

Constant rate is a strong assumption, and reality often violates it. A restaurant
gets a lunch rush; a website spikes at 9\,a.m.; earthquakes cluster after a big
one. We keep the independence and one-at-a-time structure but let the rate vary
in time, replacing the constant $\lambda$ with an \emph{intensity function}
$\lambda(t) \ge 0$.

\begin{definition}[Nonhomogeneous Poisson process]
\label{ch25:def-nonhomogeneous}
A counting process $\{N(t)\}$ with $N(0)=0$ is a \emph{nonhomogeneous Poisson
process} with intensity $\lambda(t) \ge 0$ if it has independent increments,
events occur singly, and the chance of an event in $[t, t+h)$ is $\lambda(t)h +
o(h)$. Define the \emph{mean (cumulative) intensity}
\[
m(t) = \int_0^t \lambda(u)\,\dd u.
\]
Then the count in $(a,b]$ is Poisson with mean $m(b)-m(a)$:
\[
N(b)-N(a) \sim \Pois\!\Big(\textstyle\int_a^b \lambda(u)\,\dd u\Big).
\]
\end{definition}

The mean count over a window is no longer rate $\times$ width; it is the
\emph{area under the intensity curve} over that window. Everything else is
inherited: counts in disjoint windows are still independent Poissons, only with
window-dependent means.

\begin{connectionbox}
The leap from homogeneous to nonhomogeneous is exactly the leap from
multiplication to integration. For constant rate the expected count is the
rectangle $\lambda \cdot t$; for varying rate it is the integral $\int_0^t
\lambda(u)\,\dd u$, the accumulated area under the intensity. This is the same
move that turns ``distance $=$ speed $\times$ time'' into ``distance $=
\int$ speed $\dd t$'' in the companion calculus volume. The intensity
$\lambda(t)$ is to the mean-count curve $m(t)$ exactly as velocity is to
position: $m'(t) = \lambda(t)$.
\end{connectionbox}

\begin{happeningbox}
There is a slick way to think of a nonhomogeneous process: it is a homogeneous
rate-1 process viewed through a \emph{warped clock}. Define new time
$\tau = m(t)$. Stretching the axis so that equal amounts of \emph{expected events}
occupy equal length turns the breathing rate back into a constant rate of $1$.
This ``time change'' is why so many homogeneous results carry over: you prove them
once at constant rate, then warp time. It is also exactly how you \emph{simulate}
a nonhomogeneous process --- generate a rate-1 process and unwarp, or use
\emph{thinning} (next box).
\end{happeningbox}

\begin{figure}[H]
\centering
\begin{tikzpicture}
\begin{axis}[
  width=12cm, height=5.2cm,
  domain=0:24, samples=160,
  xmin=0, xmax=24, ymin=0, ymax=14,
  xlabel={hour of day}, ylabel={intensity $\lambda(t)$ (calls/hr)},
  title={A breathing rate: two daily rushes},
  axis lines=left, title style={font=\small},
]
\addplot[name path=lam, densityred, very thick] {2 + 9*exp(-(x-12)^2/6) + 6*exp(-(x-18)^2/4)};
\addplot[name path=base, draw=none] {0};
\addplot[densityred!22] fill between[of=lam and base];
\end{axis}
\end{tikzpicture}
\caption{A nonhomogeneous intensity $\lambda(t)$ with a midday and an evening
peak. The expected number of events between two times is the \emph{area} under
this curve over that span, not a simple rate times width.}
\label{ch25:fig-nonhomog}
\end{figure}

\begin{applicationbox}[Simulation by thinning]
To simulate a nonhomogeneous process with intensity bounded by
$\lambda(t)\le \lambda^\star$, run a homogeneous rate-$\lambda^\star$ process
(easy: just add up $\Expo(\lambda^\star)$ gaps) and then \emph{keep} each
candidate event at time $t$ with probability $\lambda(t)/\lambda^\star$, throwing
the rest away. By the thinning theorem the survivors have exactly the intensity
$\lambda(t)$. This \emph{Lewis--Shedler thinning} algorithm is the standard way
to generate arrivals for nonhomogeneous queueing simulations and for fitting
point-process models in epidemiology and seismology.
\end{applicationbox}

\section{A fully worked path, end to end}
\label{ch25:worked-path}

To pin every viewpoint to one concrete computation, run them all on a single
example.

\begin{example}[Server requests]
\label{ch25:ex-server}
Requests hit a server as a Poisson process at $\lambda = 3$ per second.

\textbf{(a) Count in a window.} Requests in the next 2 seconds:
$N(2) \sim \Pois(6)$, so $\E = 6$, $\Var = 6$, and
$\P(N(2)=4) = e^{-6}6^4/4! = 1296\,e^{-6}/24 = 54\,e^{-6} \approx 0.134.$

\textbf{(b) Waiting for the next.} The next request is $\Expo(3)$ seconds away,
mean $1/3$ s. Chance you wait more than $1$ second: $e^{-3}\approx 0.0498$.

\textbf{(c) Waiting for the fifth.} $S_5\sim\Gam(5,3)$, mean $5/3$ s. Chance the
fifth request arrives within $1$ second is $\P(N(1)\ge 5)$ with
$N(1)\sim\Pois(3)$:
\[
1 - e^{-3}\!\left(1+3+\tfrac{9}{2}+\tfrac{27}{6}+\tfrac{81}{24}\right)
= 1 - e^{-3}\cdot 16.375 \approx 0.185.
\]

\textbf{(d) Thinning.} Suppose $20\%$ of requests are writes. Writes form a
Poisson process at $0.6$/s, reads at $2.4$/s, and the two streams are
independent. Writes in the next $10$ s: $\Pois(6)$.

\textbf{(e) Conditional uniformity.} Told that exactly $3$ requests arrived in
$[0,2]$, the chance all three came in the first second is, by uniformity,
$(1/2)^3 = 1/8$ --- the rate $\lambda$ has vanished entirely from the question.
\end{example}

\begin{trapbox}
A few recurring stumbles. (1) \emph{Mismatched units.} A rate is per unit time;
always convert the window to the same units before forming $\lambda t$. A rate of
$6$/hour over $10$ minutes gives mean $1$, not $6$. (2) \emph{Confusing the
exponential gap with the Gamma arrival time.} $T_i\sim\Expo(\lambda)$ is one gap;
$S_n=\sum T_i \sim\Gam(n,\lambda)$ is when the $n$-th event lands. (3)
\emph{Forgetting independence after thinning.} The two thinned streams really are
independent --- do not double-count shared randomness. (4) \emph{Assuming the gap
you observe equals the mean gap}; the inspection paradox says the gap you land in
averages $2/\lambda$.
\end{trapbox}

\begin{aestheticbox}
Step back and look at the unity of it. A handful of words --- ``events at random,
no memory, steady rate'' --- forces the Poisson count, the exponential gap, the
Gamma arrival time, the additivity of rates, the splitting into independent
streams, and the uniform scattering given the count. Nothing was bolted on; each
fact is the same assumption seen from a new angle. That is what a \emph{canonical}
object feels like in mathematics: a single idea whose every consequence is
inevitable and clean. The Poisson process is the point where ``random in time''
becomes fully computable.
\end{aestheticbox}

\section{Exercises}
\label{ch25:exercises}

\begin{exercise}[Warm-up counts]
\label{ch25:exr-counts}
Emails arrive at $\lambda = 4$ per hour as a Poisson process. Find (a) the
probability of exactly $2$ emails in the next half hour; (b) the probability of
\emph{no} emails in the next $15$ minutes; (c) the expected number and the
variance of emails in an $8$-hour workday.
\end{exercise}

\begin{exercise}[Gaps and arrival times]
\label{ch25:exr-gaps}
For the email process above, find (a) the mean and standard deviation of the time
until the next email; (b) the probability that the gap to the next email exceeds
$20$ minutes; (c) the distribution and mean of the time until the \emph{third}
email.
\end{exercise}

\begin{exercise}[The Poisson--exponential bridge]
\label{ch25:exr-bridge}
Using only the identity $\{S_n > t\} = \{N(t) < n\}$, show that
$\P(S_2 > t) = e^{-\lambda t}(1+\lambda t)$, and then differentiate to recover
the $\Gam(2,\lambda)$ density. Explain in one sentence why this identity is the
formal statement that ``counting'' and ``waiting'' are the same process.
\end{exercise}

\begin{exercise}[Superposition in practice]
\label{ch25:exr-super}
Three independent ticket machines sell at Poisson rates $5$, $8$, and $7$ per
minute. (a) What is the distribution of total sales in a $2$-minute window? (b)
Given that a sale just occurred, what is the probability it came from the second
machine? (Hint: in a tiny window the chance the event is from machine $i$ is
$\lambda_i/\sum_j\lambda_j$.)
\end{exercise}

\begin{exercise}[Thinning and independence]
\label{ch25:exr-thin}
Customers enter a shop at $\lambda = 30$ per hour; each independently makes a
purchase with probability $0.4$. (a) Find the distribution of the number of
\emph{buyers} in the next hour. (b) Find the distribution of the number of
\emph{non-buyers}. (c) If $20$ buyers walked in this hour, what does that tell you
about the number of non-buyers? Justify with the thinning theorem.
\end{exercise}

\begin{exercise}[Conditional uniformity]
\label{ch25:exr-uniform}
Exactly $5$ meteors are seen during a $4$-hour watch modeled as a Poisson
process. (a) What is the probability that exactly $2$ of them appeared in the
first hour? (b) What is the expected number that appeared in the first hour? (c)
Why does the rate $\lambda$ not appear in your answers?
\end{exercise}

\begin{exercise}[Nonhomogeneous process]
\label{ch25:exr-nonhomog}
A call center has intensity $\lambda(t) = 10 + 6\sin(\pi t/12)$ calls per hour,
$t$ in hours past midnight. (a) Write the expected number of calls between
$8$\,a.m.\ and noon as an integral and evaluate it. (b) What is the distribution
of the number of calls in that window? (c) Sketch why a constant-rate model would
mis-size staffing.
\end{exercise}

\begin{exercise}[The inspection paradox, quantitatively]
\label{ch25:exr-inspection}
Buses arrive as a Poisson process with mean gap $10$ minutes. You show up at a
uniformly random time. (a) What is the expected time until the next bus? (b) What
is the expected time since the previous bus? (c) Hence the expected length of the
gap you land in --- and explain in words why it exceeds $10$ minutes.
\end{exercise}

\begin{exercise}[Proof: memorylessness forces exponential gaps]
\label{ch25:exr-proof}
Let $G(s) = \P(T_1 > s)$ be the survival function of the first gap. Argue from
the independence and stationarity axioms that
$G(s+u) = G(s)\,G(u)$ for all $s,u\ge 0$. Then, assuming $G$ is continuous and
decreasing with $G(0)=1$, prove the only solution is $G(s) = e^{-\lambda s}$ for
some $\lambda > 0$. (This is the functional-equation heart of the whole chapter.)
\end{exercise}

\begin{exercise}[Compound Poisson, a peek ahead]
\label{ch25:exr-compound}
Claims arrive at an insurer as a Poisson process at $\lambda = 2$ per day, and
each claim is an independent amount with mean \$$1000$ and standard deviation
\$$400$. Let $X$ be the total claim amount in a $30$-day month. Using the laws of
total expectation and total variance (Chapter~15), find $\E X$ and $\Var X$.
(Hint: $\E X = \lambda T\,\E(\text{amount})$ and
$\Var X = \lambda T\,\E(\text{amount}^2)$.)
\end{exercise}
